\section{Four Powers}
\label{sec:four-powers}

A floor is not only a list of acts a system will refuse. It is an
allocation of authority, and four powers make it up: to \emph{state}
the commitment, to \emph{amend} it, to \emph{inspect} whether a
deployed system still holds it, and to \emph{enforce} it by stopping
or altering the deployment when it does not. Those four words name
those four powers from here on. Who may set any of them in motion is
a fifth question, and §\ref{sec:offices-assembled} asks it of every
office this book proposes.

Those powers need not belong to the same party, and ordinarily should
not. A laboratory may have the technical access needed to inspect the
system while lacking legitimate authority to state what the floor
protects. An outside body may have a stronger claim to state its
contents while lacking the access or leverage needed to enforce them.
A reviewer may be able to publish a finding without being able to make
anyone act on it. Calling all of these arrangements “oversight”
conceals the differences that matter.

The problem begins when one institution holds all four. A party that
can state the commitment, amend it, inspect its own compliance, and
decide whether a violation has a consequence has not bound itself. It
has described what it currently intends to do.
Section~\ref{sec:create-amend} asks who may hold the first two powers
against the institution that operates the system;
section~\ref{sec:amend-floor} asks who may exercise the second, and
section~\ref{sec:review-binding} asks what makes the last two bind.


\runin{Four powers in one hand}

An institution may have good reasons to amend a commitment, and an
unamendable rule can become dangerous or obsolete. The defect is not
amendment. It is that the party bound by the rule is the only party
entitled to amend it, the only party that inspects compliance, and
the only party that decides whether an adverse finding has a
consequence. The commitment can be sincere when it is published; what
the arrangement lacks is anything outside the institution that holds
it in place when keeping it becomes costly. Disclosure lets the public
observe a change after it operates. It gives nobody standing to
contest the change before it operates, or to require a justification.
Two cases show the powers coming apart, and which of them moved.

Google's AI Principles were published in June 2018 after employee
opposition to Project Maven, a Pentagon program that used machine
learning to analyze drone footage. Thousands of employees signed a
letter against the contract and about a dozen resigned
\autocite{castelino2018google}; the company declined to renew it
\autocite{conger2018maven}, and the principles it published excluded
weapons and surveillance \autocite{pichai2018principles}
for six and a half years, until the pledge was removed in February
2025 \autocite{google2025principles}. Run the artifact question of
§\ref{sec:antifascist-alignment} on that record: what did the demand
come back as, and who held it. It came back as a policy. An outcome
changed, the objectors had a document they could hold the company to
in public, and nobody outside the company held it. So it lasted
exactly as long as the company's willingness, and its ending had to
be announced. That is the partial case the criterion has to grade.
For six years the objection was a transfer held at the institution's
pleasure; the announcement in 2025 is what taking one back looks
like, and a recuperation would have left nothing to take back.

The removal of Timnit Gebru in December 2020 and of Margaret Mitchell
two months later is the harder case, because the power that failed to
move was inspection rather than amendment. Google employed both
researchers to examine the risks of its own systems and dismissed
them when the examination became costly
\autocite{dickey2020gebru,dickey2021mitchell}.
The defect was not a lack of expertise or the absence of an ethics
function. It was that the institution being inspected controlled the
inspectors' employment, their access to evidence, and what followed
from their findings. Inspection existed, and the power to inspect
never left the party inspected, which could point to the existence of
internal criticism as evidence of openness while keeping the power to
remove the critics whose work tested it.

For a floor meant to bind its operator, then, amendment and inspection
have to involve parties whose access, tenure, and authority do not
depend on the operator's continuing consent. The rest of this chapter
asks who those parties could be and what would make their findings
bind.
